% Root proof, 2026-11-19. Independent audit requested before promotion.
\subsection{Two-jet geometry of a genus-two theta surface in characteristic two}
\label{subsec:ppav-two-jets}
\status{Proved. The obstruction/crystal agent independently audited the complete
jet and discriminant proofs, and separately audited the supplementary dichotomy
and ordinary-member arguments.}

\begin{lemma}\label{lem:ppav-surface-dichotomy}
\status{Proved under stated hypotheses}
Over an algebraically closed field, a principally polarized abelian surface
with an effective symmetric theta divisor is either a product of two elliptic
curves with its product theta divisor, or the Jacobian of a smooth genus-two
curve with its smooth theta divisor. This assertion includes characteristic two.
\end{lemma}
\begin{proof}
An abelian variety contains no rational curve. Indeed a nonconstant map from
$\mathbf P^1$ factors through a finite number of Frobenius maps until its
differential is nonzero, whereas invariant differentials generate the
cotangent bundle of the abelian variety and $H^0(\mathbf P^1,\Omega^1)=0$.
For any integral curve $B$ on the surface, adjunction therefore gives
$B^2=2p_a(B)-2\ge0$. Also $\Theta^2=2$ and $\Theta\cdot B>0$.
Writing $\Theta=\sum m_iB_i$ shows that a reducible theta divisor must
have exactly two components, each with multiplicity one: their positive
integer intersections with $\Theta$ sum to two. A single component of
multiplicity at least two would give $m_i^2B_i^2=2$, which is impossible.
In the reducible case, evenness of $B_i^2$ and
$B_i^2+B_1\cdot B_2=1$ force $B_i^2=0$ and $B_1\cdot B_2=1$.
Each component is consequently a smooth genus-one curve, hence a translate
of an elliptic subgroup. Addition of the two elliptic subgroups is an
isogeny of degree their intersection number, namely one. It is the
required product isomorphism with its theta divisor.

If $\Theta$ is integral, $p_a(\Theta)=2$. Its normalization cannot have
genus one: a map from an elliptic curve to an abelian variety has a smooth
elliptic translate as image, which would have self-intersection zero.
Thus its normalization has genus two and arithmetic genus two, so
$\Theta$ is smooth. Its inclusion induces an isogeny
$u:\operatorname{Jac}(\Theta)\to A$, since this ample curve generates
$A$. Let $C_J$ denote the genus-two Abel--Jacobi curve in its Jacobian.
Then $u^*\Theta\cdot C_J=2$ and $C_J^2=2$. Hodge index gives
$(u^*\Theta)^2\le2$, whereas this square is $2\deg u$.
Thus $\deg u=1$ and $u$ is an isomorphism. Its restriction maps
$C_J$ onto $\Theta$, so it identifies the principal polarizations.
\end{proof}

We work in this subsection over an algebraically closed field $k$ of
characteristic two. Let $C$ be a smooth genus-two curve, choose a theta
characteristic $\kappa$ with $\kappa^{\otimes2}=\omega_C$, and put
\[
 A=\operatorname{Pic}^0(C),\qquad
 i:C\hookrightarrow A,\quad c\longmapsto\mathcal O_C(c)\otimes\kappa^{-1},
 \qquad \Theta=i(C),\qquad L=\mathcal O_A(3\Theta).
\]
Thus $\Theta$ is smooth and symmetric and $\Theta^2=2$. These Jacobian facts, including the characteristic-two
hyperelliptic description, are recorded in
\cite[Theorem~14.23 and Corollary~14.24, pp.~230--231; Theorem~14.29, p.~233;
Example~14.34, pp.~236--237]{EdixhovenVanDerGeerMoonenAV}.
Write $\Theta_a=\Theta+a$ and
$P_a=t_a^*\mathcal O_A(\Theta)\otimes\mathcal O_A(-\Theta)$.
The principal polarization identifies $a\mapsto P_a$ with an isomorphism
$A\to\operatorname{Pic}^0(A)$. Additions and integer multiples of points
below are in this group law. Put $W=H^0(A,L)$, so $\dim W=9$.
For a point $P\in A(k)$, a first jet means the value in
$L_P\otimes\mathcal O_{A,P}/\mathfrak m_P^2$, of dimension three;
a second jet means the analogous quotient by $\mathfrak m_P^3$, of
dimension six. All jet maps in this subsection are $k$-linear evaluation maps.

\begin{lemma}\label{lem:ppav-theta-preliminaries}
\status{Proved under stated hypotheses}
Every line bundle $\mathcal O_A(2\Theta)\otimes P_a$ is globally generated
and has no positive-degree cohomology. The system $|2\Theta|$ separates
two distinct points $x,y$ whenever $y\ne -x$.
Moreover, under $\operatorname{Pic}^0(C)=A$,
\[
                     i^*P_a=-a.
\]
The Gauss map of $\Theta$ to the projective line of cotangent directions
is the canonical degree-two map of $C$.
\end{lemma}

\begin{proof}
The theorem of the square gives
$\Theta_c+\Theta_{-c}\in|2\Theta|$.
For a specified point $x$, choose $c$ outside the two curves
$\Theta+x$ and $\Theta-x$; the resulting section is nonzero at $x$.
If $y\ne x,-x$, choose $c\in\Theta+x$ outside
$(\Theta+y)\cup(\Theta-y)$. These are distinct irreducible translates:
the translation stabilizer of $\Theta$ is contained in the kernel of
its principal polarization and is trivial. Their intersections with
$\Theta+x$ are finite, so this choice exists. The resulting section
vanishes at $x$ and is nonzero at $y$. Together with a section nonzero
at both points this proves separation. Every twist by $P_a$ is a
translate of $\mathcal O(2\Theta)$, since multiplication by two on
$A(k)$ is surjective, so it is globally generated as well.

For completeness the required ample vanishing does not use
characteristic-zero Kodaira vanishing. If $n$ is odd, the finite etale
map $[n]:A\to A$ has invertible degree $n^4$, and its trace splits the
unit map of structure sheaves. Thus
$H^j(A,T)\hookrightarrow H^j(A,[n]^*T)$ for a line bundle $T$.
For $T=\mathcal O(2\Theta)$, symmetry gives
$[n]^*T\simeq\mathcal O(2n^2\Theta)$; Serre vanishing for arbitrarily
large odd $n$ proves the assertion for $j>0$. Translation proves it
for the twists. The same argument and Riemann--Roch give
$h^0(3\Theta)=(3\Theta)^2/2=9$.

Let $\iota$ be the hyperelliptic involution of $C$; then
$i(\iota c)=-i(c)$ and $c+\iota c$ is canonical. For a generic
$a=i(r)-i(s)$, the intersection $\Theta\cap\Theta_a$ consists of
$i(r)$ and $i(\iota s)$. Indeed these are intersection points, they
are distinct generically, and the intersection has length two.
The difference map $\Theta\times\Theta\to A$ is surjective, since
two ample translates intersect. Restricting $\mathcal O(\Theta_a)$
to $C$ therefore gives generically
$\mathcal O_C(r+\iota s)$. Since
$\mathcal O(\Theta_a)=\mathcal O(\Theta)\otimes P_{-a}$ and
$i^*\mathcal O(\Theta)=\omega_C=\mathcal O_C(s+\iota s)$,
we obtain $i^*P_{-a}=a$. Equality extends from this dense set to the
morphisms $A\to\operatorname{Pic}^0(C)$, proving the formula for all $a$.
Finally the restriction of invariant one-forms on the Jacobian to
$C$ is the defining isomorphism
$H^0(A,\Omega_A^1)\simeq H^0(C,\omega_C)$.
The conormal direction of $\Theta$ is consequently its canonical
map. This map has degree two and is finite and surjective onto
$\mathbf P^1$. No assertion about a characteristic-two Kummer
quotient is needed.
\end{proof}

\begin{lemma}\label{lem:ppav-three-distinct-values}
\status{Proved under stated hypotheses}
In the smooth-theta case, any three distinct points of $A$ impose three
independent value conditions on $H^0(A,3\Theta)$.
\end{lemma}
\begin{proof}
It suffices to separate the value at $P_3$ from zero values at
$P_1,P_2$. Choose $a\in\Theta+P_1$ and $b\in\Theta+P_2$.
Require $a,b\notin\Theta+P_3$ and $P_3+a+b\notin\Theta$.
The first two conditions define dense open subsets of their respective
irreducible curves, because the three points are distinct and a theta
divisor has trivial translation stabilizer. The third defines a
nonempty open subset of their product: the sum of these two translates
of $\Theta$ is all of $A$. These three opens have nonempty intersection.
The product of theta sections for $\Theta_a,\Theta_b,\Theta_{-a-b}$
then vanishes at $P_1,P_2$ and not at $P_3$. Repeating with each
point in the third position proves the assertion.
\end{proof}

\begin{theorem}\label{thm:ppav-three-theta-jets}
\status{Proved under stated hypotheses}
With the data just specified, the following assertions hold.
\begin{enumerate}
\item $L$ separates points and first jets, and it separates second jets
at every point of $A$.
\item For $P\ne Q$, the map
\[
 W\longrightarrow
 L_P\otimes\mathcal O_{A,P}/\mathfrak m_P^2\ \oplus\
 L_Q\otimes\mathcal O_{A,Q}/\mathfrak m_Q^2
\]
has rank at least four. It has rank six outside the dimension-two locus
\[
 R=\{(P,Q):2P+Q\in\Theta,\ P+2Q\in\Theta\}.
\]
Its rank is at least five outside the dimension-at-most-one locus
\[
 R_4=\bigcup_{s\in A[3](k)}
       \{(P,s-P):P+s\in\Theta\}.
\]
The rank-four locus is contained in $R_4$, whether or not the rank
actually drops to four at each listed pair.
\item If $A$ is ordinary, then for every fixed $P$ the fibre $R_P$
is finite. In particular two first jets have full rank six for all
but finitely many $Q$ on any specified integral ample curve.
\end{enumerate}
\end{theorem}

\begin{proof}
For first jets at $P$, choose a theta translate through $P$ and
multiply its equation by a section of the remaining twist of
$2\Theta$ nonzero at $P$. The possible linear terms span both
cotangent directions because the Gauss map of $\Theta$ is
surjective. A product nonzero at $P$ supplies the value. Point
separation follows by choosing a theta translate through one point
and not the other, and a nonzero section of the remaining twist
at the latter point. Thus $L$ is very ample as well.

For second jets, fix $P$ and take $p,q\in\Theta$. Put
$a=P-p$, $b=P-q$, and $c=-a-b$. The product of the three theta
sections for $\Theta_a,\Theta_b,\Theta_c$ belongs to $|3\Theta|$.
Its first two factors vanish at $P$; the third is nonzero precisely
when $3P-p-q\notin\Theta$. This condition defines a nonempty open
subset of $\Theta\times\Theta$, because the addition map to $A$
is surjective. The product of the two Gauss maps is finite and
surjective onto $\mathbf P^1\times\mathbf P^1$. The image of that
open is dense, and the associated products of two linear forms
span all of $\operatorname{Sym}^2T_P^*A$: multiplication of binary
linear forms has dense image in the projective space of binary
quadratic forms, also in characteristic two. The resulting
quadratic terms therefore span
$\mathfrak m_P^2/\mathfrak m_P^3$. Combined with first-jet
surjectivity this proves second-jet surjectivity.

Now write $s=P+Q$. The theta translates through both $P,Q$ are
parametrized by the length-two scheme
\[
                      Z=(\Theta+P)\cap(\Theta+Q).
\]
It is preserved by $a\mapsto s-a$. A fixed point of this involution
has $2a=s$; the corresponding points $P-a,Q-a$ on $\Theta$ are
hyperelliptic conjugates and have equal Gauss direction, so its
intersection multiplicity is at least two. Consequently two
distinct support points of $Z$ are exchanged, and a single support
point is fixed.

Fix one such parameter $a$, and write $p=P-a$, $q=Q-a$ as points
of $C$. The restriction exact sequence is
\[
 0\longrightarrow H^0(2\Theta\otimes P_a)
 \longrightarrow W\longrightarrow H^0(\Theta_a,L|_{\Theta_a})
 \longrightarrow0;
\]
the surjectivity follows from Lemma~\ref{lem:ppav-theta-preliminaries}.
After translation to $C$, its last line bundle is
$M_a=\omega_C^{\otimes3}\otimes(-3a)$, of degree six.
The four along-curve jet conditions have rank at least three.
Indeed their cokernel is $H^1(C,M_a(-2p-2q))$, which has dimension
at most one for this degree-two line bundle. It is one exactly
when $M_a(-2p-2q)=\omega_C$. Since
$\mathcal O_C(p+q)\otimes\omega_C^{-1}$ has class $s-2a$,
this happens exactly when
\[
                              a=2s.
\]
The sections vanishing on $\Theta_a$ supply normal derivatives of
rank at least one. They supply rank two if $a\ne -s$: choose
$b\in A(k)$ with $2b=a$. Their evaluations are those of
$2\Theta$ at $P+b,Q+b$, which are distinct and are not opposite
unless $a=-s$. Lemma~\ref{lem:ppav-theta-preliminaries} applies.
Projection of the surface jet space to the along-curve jet space
has the two normal directions as its kernel. Hence the full rank
is at least the sum of the along-curve rank and the supplied normal
rank. It is therefore at least four always, six if a parameter
$a$ avoids both $-s,2s$, and five unless those two bad parameters
coincide.

If every support point of $Z$ is among $-s,2s$, then $-s$ is a
support point. For two support points this is immediate. For one
support point equal to $2s$, its being fixed gives $4s=s$, so
$3s=0$ and $2s=-s$. Thus rank failure from six implies that
$\Theta_{-s}$ contains both $P,Q$, which is exactly the displayed
condition defining $R$. Rank below five requires $3s=0$, and then
these two containment conditions reduce to $P+s\in\Theta$.
The map $(P,Q)\mapsto(2P+Q,P+2Q)$ is an etale isogeny of
$A\times A$ (its determinant is three). Its inverse image of
$\Theta\times\Theta$ has dimension two. Since $A[3](k)$ has
$3^4=81$ elements, $R_4$ has dimension at most one.

For the final assertion suppose a fibre $R_P$ were positive
dimensional. Its first condition places it on the irreducible curve
$\Theta-2P$; the second would then force
\[
                           [2]\Theta=\Theta+3P.
\]
When $A$ is ordinary, write $[2]=V\circ F_{A/k}$, where $V$ is
etale. The restriction of $F_{A/k}$ to $\Theta$ has degree two.
The restriction of $V$ to $\Theta^{(2)}$ has generic degree one:
any larger degree would make that curve invariant under a nonzero
element of the constant kernel of $V$, whereas a principal theta
divisor has trivial translation stabilizer. Thus $[2]|_\Theta$
has degree two. The asserted image equality contradicts the
projection formula: its left interpretation gives
$([2]_*\Theta)\cdot\Theta=2\Theta^2=4$, while the projection
formula gives $\Theta\cdot[2]^*\Theta=4\Theta^2=8$.
This proves finiteness and the last assertion.
\end{proof}

\begin{corollary}\label{cor:ppav-discriminant-one-node}
\status{Proved under stated hypotheses}
The discriminant of singular divisors in $|3\Theta|$ is irreducible
of dimension seven. Its general member is integral and has exactly
one singularity, an ordinary node. These assertions hold for every
$C$ in the stated characteristic; ordinary $A$ is not needed here.
\end{corollary}

\begin{proof}
The incidence of a section singular at a specified point is a
$\mathbf P^5$-bundle over $A$, hence is irreducible of dimension
seven. Second-jet surjectivity makes an ordinary node a nonempty
open condition on this incidence: the mixed coefficient in a
binary quadratic form is nonzero. The incidence of two distinct
singular points has dimension at most six. Indeed its rank-six
stratum has dimension at most $4+(8-6)=6$, its rank-five stratum
at most $2+(8-5)=5$, and its rank-four stratum at most
$1+(8-4)=5$. A general point of the single-point incidence thus
has exactly one singular point and it is an ordinary node.
Projection to $|3\Theta|$ is generically one-to-one, proving the
dimension and irreducibility assertions.

Such a divisor is reduced, since a multiple component would give
a positive-dimensional singular locus. It cannot be reducible:
if it splits as nonzero effective divisors $B+D\sim3\Theta$,
put $t=\Theta\cdot B$, so $1\le t\le5$. Hodge index gives
$B^2\le t^2/2$, and hence
\[
 B\cdot D=3t-B^2\ge 3t-t^2/2\ge5/2.
\]
Its intersection multiplicity is therefore at least three, whereas
two components of a curve with a unique ordinary node can have
total intersection multiplicity at most one. This contradiction
proves integrality.
\end{proof}

\begin{lemma}\label{lem:ppav-ordinary-members}
\status{Proved under stated hypotheses}
If the genus-two Jacobian $A$ above is ordinary, $|3\Theta|$ contains
an ordinary connected nodal divisor with three smooth theta components.
The general discriminant member is ordinary. For every $d\ge1$, a general
smooth member of $|3d\Theta|$ is ordinary, of genus $9d^2+1$.
\end{lemma}
\begin{proof}
Take the three translates $\Theta_a,\Theta_b,\Theta_{-a-b}$.
Their sum belongs to $|3\Theta|$ by the theorem of the square.
They are distinct and pairwise transverse for general $(a,b)$.
To see transversality without a generic-smoothness assumption in positive
characteristic, the difference map $\Theta^2\to A$ is generically etale:
its differential is invertible when the two Gauss directions differ.
Its critical values thus form a proper closed subset of $A$.
Each of the differences $b-a,-2a-b,-a-2b$ varies surjectively in $A$,
so a general pair avoids all three exceptional subsets.
Triple intersections are absent for a general pair as well. A triple
intersection corresponds to $x,p,q,r$ with $p,q,r\in\Theta$ and
$3x=p+q+r$, followed by $a=x-p,b=x-q$. Since $[3]$ is finite,
this incidence has dimension three, below the four-dimensional
parameter space of $(a,b)$.

The three components are ordinary because their Jacobian is the ordinary
surface $A$. The normalization sequence for a nodal curve expresses its
$H^1(\mathcal O)$ as an extension of the component cohomologies by its
graph cohomology. Frobenius is bijective on both terms, so this union
is ordinary. Its ordinary point belongs to the irreducible discriminant
of Corollary~\ref{cor:ppav-discriminant-one-node}, whose ordinary open
is therefore dense. Ordinarity is also a nonempty open condition in
the entire linear system. Its intersection with the nonempty smooth
open supplies smooth ordinary members of $|3\Theta|$.

Finally take $d$ general such members, meeting transversely without
triple points. Their nodal union lies in $|3d\Theta|$ and is ordinary
by the same normalization sequence. The positive-degree ample vanishing
used above makes $h^1(\mathcal O_Y)=9d^2+1$ constant for this family
of effective divisors. Invertibility of its coherent Frobenius is
therefore open, and intersecting again with the smooth open gives
the asserted general ordinary curve and genus.
\end{proof}

\begin{remark}[Source scope]\label{auto:A_ppav_two_jets:0557c5a87cfe}
\status{Proved under stated hypotheses}

The source search found Bauer--Szemberg's higher-jet statements
\cite[Theorems~1 and~2, conventions p.~2]{BauerSzemberg1997}
and Alvarado--Pareschi's recent three-theta second-jet theorem
\cite[Theorem~D; Section~2 conventions]{AlvaradoPareschi2026},
but their stated conventions are in characteristic zero. We do
not import either statement in characteristic two. The proof
above uses the theta theorem of the square, adjunction on a
genus-two curve, the elementary product sections of $2\Theta$,
and the explicit incidence calculation. The fixed-point finiteness
assertion explicitly requires ordinary $A$.
\end{remark}
